\section*{Hoofdstuk \ref{ch:b3:seismology-imaging} --- Seismologie en aardbeeldvorming}

\begin{solution}{exo:b3:seismology-imaging:1}
(a) $\sin i_0 = pv_0/R = 500\times5.80/6371 = 0.4552$, dus $i_0 = \ang{27.1}$.
(b) Bij \qty{771}{km} is $r = \qty{5600}{km}$ en $r/v = 5600/11.066 =
\qty{506.1}{s.rad^{-1}}$; bij \qty{871}{km} is $r = \qty{5500}{km}$ en $r/v =
5500/11.245 = \qty{489.1}{s.rad^{-1}}$. De straal keert om waar $r/v = p = 500$:
$z_t = 771 + 100\times(506.1 - 500)/(506.1 - 489.1) = 771 + 36 =
\qty{807}{km}$.
\end{solution}

\begin{solution}{exo:b3:seismology-imaging:2}
(a) $(455.5 - 369.6)/10 = \qty{8.59}{s}$ per graad, en
$8.59\times180/\pi = \qty{492.2}{s.rad^{-1}}$. (b) Eén boograad is
$6371\pi/180 = \qty{111.19}{km}$, dus het golffront doorkruist het oppervlak met
$111.19/8.59 = \qty{12.94}{km.s^{-1}}$, veel sneller dan elk gesteente aan het
oppervlak: de straal komt steil aan. (c) Op het keerpunt geldt $v = r_t/p =
5450/492.2 = \qty{11.07}{km.s^{-1}}$, een P-snelheid in de onderste mantel.
\end{solution}

\begin{solution}{exo:b3:seismology-imaging:3}
(a) $\mu_1 = \rho_1\beta_1^2 = 2700\times3500^2 = \qty{3.31e10}{Pa}$
(\qty{33.1}{GPa}); $\mu_2 = 3300\times4500^2 = \qty{6.68e10}{Pa}$
(\qty{66.8}{GPa}). (b) $\lambda = cT = 3.79\times20 = \qty{75.8}{km}$, 2.17
maal de laagdikte: de golf reikt ver de halfruimte in, en dat is
waarom zijn snelheid al boven $\beta_1$ ligt. (c) $4000/3.384 = \qty{1182}{s}$ en
$4000/4.333 = \qty{923}{s}$: de 20-s-golven komen \qty{259}{s}, of
\qty{4.3}{min}, later aan dan de 100-s-golven.
\end{solution}

\begin{solution}{exo:b3:seismology-imaging:4}
(a) ${}_0S_2$: $1/(53.89\times60) = \qty{3.093e-4}{Hz} = \qty{0.3093}{mHz}$;
${}_0T_2$: $1/(44.01\times60) = \qty{0.3787}{mHz}$. (b) $1440/53.89 = 26.7$
en $1440/44.01 = 32.7$ trillingen per dag. (c) Het verschil is
$\qty{0.0694}{mHz}$, waarvan de inverse \qty{14400}{s} is: minstens \qty{4.0}{h} aan
registratie, in de praktijk veel meer, omdat de pieken door demping worden verbreed en
gesplitst.
\end{solution}

\begin{solution}{exo:b3:seismology-imaging:5}
(a) De koorde tussen twee punten van de bol, gescheiden door de hoek $\Delta$
in het middelpunt heeft de lengte $2R\sin(\Delta/2)$, afgelegd met $v$.
(b) Bij \ang{30}: $2\times6371\times\sin\ang{15}/10.0 = \qty{329.8}{s}$, tegenover
\qty{369.6}{s}; bij \ang{90}: $2\times6371\times\sin\ang{45}/10.0 =
\qty{901.0}{s}$, tegenover \qty{779.7}{s}. (c) De straal naar \ang{30} keert om in de
bovenmantel, waar de werkelijke snelheid ongeveer 8 tot \qty{10}{km.s^{-1}} is, lager dan de
uniforme \qty{10}{km.s^{-1}}; de straal naar \ang{90} doorloopt het grootste deel van zijn weg in de
onderste mantel, waar de snelheid 11 tot \qty{13.7}{km.s^{-1}} is, en buigt er naartoe af, zodat de
werkelijke Aarde is sneller.
\end{solution}

\begin{solution}{exo:b3:seismology-imaging:6}
(a) Na kwadrateren: $(2 - \xi)^4 = 16(1 - \xi/3)(1 - \xi)$. Werk uit:
$16 - 32\xi + 24\xi^2 - 8\xi^3 + \xi^4 = 16 - \tfrac{64}{3}\xi + \tfrac{16}{3}\xi^2$,
dus $\xi^4 - 8\xi^3 + \tfrac{56}{3}\xi^2 - \tfrac{32}{3}\xi = 0$, dat is
$\xi\,(3\xi^3 - 24\xi^2 + 56\xi - 32) = 0$, en
$3\xi^3 - 24\xi^2 + 56\xi - 32 = (\xi - 4)(3\xi^2 - 12\xi + 8)$. De wortel
$\xi = 0$ is geen golf. (b) De wortels van de kwadratische vergelijking zijn $2 \pm 2/\sqrt3$,
dat is 3.155 en 0.8453; alleen $\xi = 0.8453$ ligt in $(0, 1)$ (en $\xi = 4$
niet), en hij voldoet aan de ongekwadrateerde vergelijking, want beide zijden zijn $4/3$.
Dus $c/\beta = \sqrt{0.8453} = 0.9194$. (c) $c = 0.9194\times3.46 =
\qty{3.181}{km.s^{-1}}$; over \qty{1000}{km} is het verschil $314.4 - 289.0 = \qty{25.3}{s}$
na de S-golf.
\end{solution}

\begin{solution}{exo:b3:seismology-imaging:7}
$\mathsf G$ heeft de rijen $(1, x_i)$ en $\mathsf C_d^{-1} = \operatorname{diag}
(100, 100, 25, 25)$ (\unit{s^{-2}}). De normaalvergelijkingen zijn
\[
  \begin{pmatrix} 250 & 23750 \\ 23750 & \num{2812500} \end{pmatrix}
  \begin{pmatrix} a \\ b \end{pmatrix}
  = \begin{pmatrix} 3665 \\ 418125 \end{pmatrix},
\]
met determinant \num{1.390625e8}. Dus $b = (250\times418125 -
23750\times3665)/\num{1.390625e8} = \qty{0.12575}{s.km^{-1}}$ en
$a = (2812500\times3665 - 23750\times418125)/\num{1.390625e8} =
\qty{2.713}{s}$. De covariantie is $\mathsf N^{-1}$: $\sigma_b =
\sqrt{250/\num{1.390625e8}} = \qty{0.00134}{s.km^{-1}}$, $\sigma_a =
\sqrt{2812500/\num{1.390625e8}} = \qty{0.142}{s}$. De snelheid is
$1/b = \qty{7.95}{km.s^{-1}}$, met de fout $\sigma_b/b^2 =
\qty{0.08}{km.s^{-1}}$: een snelheid van de bovenste mantel. De residuen zijn $-0.001$,
$0.011$, $-0.076$ en $\qty{0.036}{s}$, en $\chi^2 = 0.19$ voor $4 - 2 = 2$
vrijheidsgraden: kleiner dan verwacht, dus de opgegeven fouten zijn eerder
ruim; de fit is aanvaardbaar.
\end{solution}

\begin{solution}{exo:b3:seismology-imaging:8}
(a) $\mathsf G = (30\;\;40)$ is een enkele rij: $\mathsf G^{\mathsf T}\mathsf G$ is
$2\times2$ met rang één, dus singulier; elk model met $30m_1 + 40m_2 = 0$ kan
worden toegevoegd zonder het gegeven te veranderen. (b) $\mathsf G\mathsf G^{\mathsf T} =
900 + 1600 = 2500$, dus $\mathbf m = (30, 40)\times0.50/2500 =
(0.006, 0.008)\,\unit{s.km^{-1}}$. (c) $\mathsf R = \frac{1}{2500}
\begin{pmatrix} 900 & 1200 \\ 1200 & 1600 \end{pmatrix} =
\begin{pmatrix} 0.36 & 0.48 \\ 0.48 & 0.64 \end{pmatrix}$, met spoor 1:
één gegeven bepaalt één combinatie van parameters. (d) $\mathsf R\,(0.01, 0) =
(0.0036, 0.0048)\,\unit{s.km^{-1}}$: de anomalie is uitgespreid langs de straal, meer
in de cel waar de straal langer is.
\end{solution}

\begin{solution}{exo:b3:seismology-imaging:9}
(a) Gelineariseerd: $\delta t = -(\delta v/v)L/v = -0.03\times200/9.0 =
\qty{-0.667}{s}$; exact, $L/(1.03v) - L/v = 200/9.27 - 200/9.0 =
\qty{-0.647}{s}$. (b) Het pad in de slab is $100/\cos\ang{40} =
\qty{130.5}{km}$: $-0.03\times130.5/9.0 = \qty{-0.435}{s}$ gelineariseerd en
$\qty{-0.422}{s}$ exact. (c) De exacte waarde is de gelineariseerde gedeeld door
$1.03$: een fout van \qty{3}{\percent}, van de orde van $\delta v/v$ zelf, zoals
verwacht van een eerste-orde theorie.
\end{solution}

\begin{solution}{exo:b3:seismology-imaging:10}
(a) De boventoon bestaat zolang zijn deel in de halfruimte vervalt, $c < \beta_2$; bij zijn
afsnijpunt $c = \beta_2$, dus $q_2 = 0$ en $\tan(kHq_1) = 0$, met $kHq_1$ in
$[n\pi, n\pi + \pi/2)$: $kHq_1 = n\pi$. (b) Met $k = \omega/\beta_2$ en
$q_1 = \sqrt{\beta_2^2/\beta_1^2 - 1} = \sqrt{\beta_2^2 - \beta_1^2}/\beta_1$,
$\omega_nH\sqrt{\beta_2^2 - \beta_1^2}/(\beta_1\beta_2) = n\pi$, wat de
formule is. (c) $\sqrt{4.5^2 - 3.5^2} = \sqrt8 = 2.828$; $\omega_1 =
\pi\times3.5\times4.5/(35\times2.828) = \qty{0.4998}{rad.s^{-1}}$, een periode van
$2\pi/0.4998 = \qty{12.57}{s}$: bij langere perioden is alleen de grondmode
aanwezig. (d) Voor $n = 0$ is de afsnijfrequentie nul: de grondmode
bestaat bij alle perioden. Bij korte perioden $kH \to \infty$ terwijl $kHq_1$ blijft
onder $\pi/2$, dus $q_1 \to 0$ en $c \to \beta_1$.
\end{solution}

\begin{solution}{exo:b3:seismology-imaging:11}
(a) $\mathsf G^{\mathsf T}(\mathsf G\mathsf G^{\mathsf T} + \varepsilon^2\mathsf I)
= (\mathsf G^{\mathsf T}\mathsf G + \varepsilon^2\mathsf I)\mathsf G^{\mathsf T}$;
vermenigvuldig links met $(\mathsf G^{\mathsf T}\mathsf G + \varepsilon^2
\mathsf I)^{-1}$ en rechts met $(\mathsf G\mathsf G^{\mathsf T} +
\varepsilon^2\mathsf I)^{-1}$; beide zijn inverteerbaar voor $\varepsilon > 0$.
(b) Hier is $\mathsf G\mathsf G^{\mathsf T} = \|\mathbf g\|^2 = 2500$, een getal,
dus $\mathsf R_\varepsilon = \mathsf G^{\mathsf T}\mathsf G/(2500 + \varepsilon^2)$,
met spoor $2500/(2500 + \varepsilon^2)$. (c) Het voorspelde gegeven is
$\mathbf g\cdot\mathbf m_\varepsilon = 0.50\times2500/(2500 + \varepsilon^2)$
en het residu $0.50\,\varepsilon^2/(2500 + \varepsilon^2)$. Voor
$\varepsilon^2 = 0$: spoor 1, residu 0; voor 2500: spoor 0.5, residu
\qty{0.25}{s}; voor \num{10000}: spoor 0.2, residu \qty{0.40}{s}. Meer demping
geeft een kleiner model, een slechtere resolutie en een grotere misfit.
\end{solution}

\begin{solution}{exo:b3:seismology-imaging:12}
(a) Uit opgave 5 volgt $T = (2R/v)\sin(\Delta/2)$, dus $\mathrm dT/\mathrm d\Delta
= (R/v)\cos(\Delta/2)$. Een koorde die onder hoek $\Delta$ in het middelpunt staat, is
de basis van een gelijkbenige driehoek met tophoek $\Delta$; zij snijdt de straal
naar de bron dus onder $i_0 = (\pi - \Delta)/2$, en $p = R\sin i_0/v =
(R/v)\cos(\Delta/2)$: hetzelfde. (b) Het keerpunt is het midden van de
koorde, op afstand $R\cos(\Delta_1/2)$ van het middelpunt; gelijkwaardig
$r_1 = p_1v$. (c) Hier is $p(\Delta)/p_1 = \cos(\Delta/2)/\cos(\Delta_1/2)$, en
de formule geeft $r_1 = R\exp\bigl(-\ln(1/\cos(\Delta_1/2))\bigr) =
R\cos(\Delta_1/2)$, zoals in (b); de snelheid is $r_1/p_1 = R\cos(\Delta_1/2)\big/
\bigl((R/v)\cos(\Delta_1/2)\bigr) = v$.
\end{solution}

\begin{solution}{pb:b3:seismology-imaging:1}
\textbf{1.} $i_0 = \arcsin 0.8 = \ang{53.13}$ en $p = 0.8/5.0 =
\qty{0.16}{s.km^{-1}}$.

\textbf{2.} De straal keert om waar $v = 1/p = \qty{6.25}{km.s^{-1}}$, dat is bij
$z = (6.25 - 5.0)/0.05 = \qty{25}{km}$.

\textbf{3.} Door $\sin i = pv$ langs de straal te differentiëren volgt $\cos i\,\mathrm di/
\mathrm ds = pg\,\mathrm dz/\mathrm ds = pg\cos i$, dus $\mathrm di/\mathrm ds =
pg$: de richting draait met een constante snelheid, de kromming van een cirkel met straal
$1/(pg) = 1/(0.16\times0.05) = \qty{125}{km}$. In het keerpunt is de straal
horizontaal op \qty{25}{km} diepte en ligt het middelpunt loodrecht daarboven op
\qty{125}{km}: bij $z = \qty{-100}{km}$, dat is $v_0/g$ boven het oppervlak, waar de
geëxtrapoleerde snelheid nul zou worden.

\textbf{4.} In de bron staat de straal van de cirkel loodrecht op de straal, dus maakt
die de hoek $i_0$ met de horizontaal, en ligt het middelpunt horizontaal
$125\cos i_0$ verder; uit symmetrie volgt
$X = 2\times125\times\cos i_0 = 2\times125\times0.6 = \qty{150}{km}$, of
$X = 2\sqrt{1 - p^2v_0^2}/(pg)$.

\textbf{5.} $\mathrm dT = \mathrm ds/v$ met $\mathrm ds = \mathrm di/(pg)$ en
$v = \sin i/p$: $\mathrm dT = \mathrm di/(g\sin i)$. Dan
\[
  T = \frac2g\int_{i_0}^{\pi/2}\frac{\mathrm di}{\sin i} = \frac2g\ln\cot\frac{i_0}2 =
  \frac2g\operatorname{arccosh}\frac1{pv_0} .
\]
Hier is $\cot(\ang{26.57}) = 2$, dus
$T = 40\ln2 = \qty{27.73}{s}$.

\textbf{6.} Voor $p = 0.15$: $pv_0 = 0.75$, $\sqrt{1 - 0.5625} = 0.6614$,
$X = 2\times0.6614/0.0075 = \qty{176.4}{km}$ en $T = 40\operatorname{arccosh}
(1.3333) = 40\times0.7954 = \qty{31.81}{s}$. Dan is $\Delta T/\Delta X =
4.08/26.4 = \qty{0.155}{s.km^{-1}}$, tussen 0.15 en 0.16, zoals
$\mathrm dT/\mathrm dX = p$ vereist.

\textbf{7.} $\mu_1 = 2700\times3500^2 = \qty{33.1}{GPa}$ en
$\mu_2 = 3300\times4500^2 = \qty{66.8}{GPa}$.

\textbf{8.} $k = 2\pi/(cT) = 2\pi/(3.99\times25) = \qty{0.06299}{rad.km^{-1}}$;
$q_1 = \sqrt{(3.99/3.5)^2 - 1} = 0.5474$; $q_2 = \sqrt{1 - (3.99/4.5)^2} = 0.4624$.

\textbf{9.} $\tan(kHq_1) = 66.825\times0.4624/(33.075\times0.5474) = 1.707$; de
grondmode neemt $kHq_1 = \arctan1.707 = 1.041$, dus
$H = 1.041/(0.06299\times0.5474) = \qty{30.2}{km}$.

\textbf{10.} Bij \qty{40}{s}: $k = 2\pi/170 = 0.03696$, $q_1 = 0.6888$,
$q_2 = 0.3287$, $\tan(kHq_1) = 0.964$, $kHq_1 = 0.7671$ en
$H = 0.7671/(0.03696\times0.6888) = \qty{30.1}{km}$. De twee komen overeen binnen de
nauwkeurigheid van de fasesnelheden: de korst is ongeveer \qty{30}{km} dik.

\textbf{11.} $\omega = 0.2513$ en $\qty{0.1571}{rad.s^{-1}}$, $k = 0.06299$ en
$\qty{0.03696}{rad.km^{-1}}$: $U \approx 0.0942/0.02603 = \qty{3.62}{km.s^{-1}}$,
kleiner dan beide fasesnelheden, zoals het moet zijn wanneer $c$ toeneemt met de
periode.

\textbf{12.} $\lambda = 3.99\times25 = \qty{99.8}{km}$ en $4.25\times40 =
\qty{170}{km}$. De grondmode bemonstert tot een fractie van zijn golflengte: de golf
van 25~s ligt grotendeels in de korst, de golf van 40~s reikt ver de mantel in, zodat
zijn snelheid dichter bij $\beta_2$ ligt. Twee perioden die dezelfde $H$ opleveren, is
een controle dat het tweelagenmodel voldoende is.

\textbf{13.} Met de kolommen voor cellen 1 tot 4 en de rijen A tot F, in eenheden van
\qty{1000}{km} ($\sqrt2\times100/1000 = 0.1414$):
\[
  \mathsf G = \begin{pmatrix}
  0.1 & 0.1 & 0 & 0 \\ 0 & 0 & 0.1 & 0.1 \\ 0.1 & 0 & 0.1 & 0 \\
  0 & 0.1 & 0 & 0.1 \\ 0.1414 & 0 & 0 & 0.1414 \\ 0 & 0.1414 & 0.1414 & 0
  \end{pmatrix}.
\]

\textbf{14.} $\mathsf C_d^{-1} = \operatorname{diag}(400, 400, 400, 400, 100,
100)$. Een rij- of kolomstraal draagt $0.1^2\times400 = 4$ bij, een diagonale straal
$0.1414^2\times100 = 2$. Cel 1 wordt gekruist door A, C (elk 4) en E (2): 10. Cellen
1 en 2 delen straal A: 4; cellen 1 en 3 delen C: 4; cellen 1 en 4 delen E: 2; en
zo verder door symmetrie, wat de vier gegeven rijen oplevert.

\textbf{15.} Een rij- of kolomstraal levert $0.1\times400\,\delta t = 40\,
\delta t$ aan elk van zijn cellen, een diagonale straal $0.1414\times100\,\delta t =
14.14\,\delta t$. Cel 1: $40(0.16 - 0.36) + 14.14\times(-0.15) = -10.12$; cel
2: $40(0.16 + 0.11) + 14.14\times(-0.27) = 6.98$; cel 3: $40(-0.48 - 0.36) +
14.14\times(-0.27) = -37.42$; cel 4: $40(-0.48 + 0.11) + 14.14\times(-0.15) =
-16.92$.

\textbf{16.} Voor $(1,1,1,1)$ somt elke rij van $\mathsf N$ op tot 20. Voor
$(1,-1,-1,1)$ geeft rij 1 $10 - 4 - 4 + 2 = 4$, en de andere rijen hetzelfde
patroon. Voor $(1,1,-1,-1)$ geldt $10 + 4 - 4 - 2 = 8$, en voor $(1,-1,1,-1)$
$10 - 4 + 4 - 2 = 8$. De vectoren staan loodrecht op elkaar en hebben norm 2, dus
$\mathsf N^{-1} = \sum_k \mathbf u_k\mathbf u_k^{\mathsf T}/(4\lambda_k)$; het
$(1,1)$-element is $\tfrac14(\tfrac1{20} + \tfrac14 + \tfrac18 + \tfrac18) =
\tfrac{11}{80}$, het $(1,2)$-element $\tfrac14(\tfrac1{20} - \tfrac14 + \tfrac18
- \tfrac18) = -\tfrac4{80}$, het $(1,4)$-element $\tfrac14(\tfrac1{20} +
\tfrac14 - \tfrac18 - \tfrac18) = \tfrac1{80}$, en zo verder.

\textbf{17.} $\hat{\mathbf m} = \mathsf N^{-1}\mathbf b$:
$\hat m_1 = (11\times(-10.121) - 4\times6.982 - 4\times(-37.418) - 16.921)/80
= -0.081$; op dezelfde manier $\hat m_2 = 1.844$, $\hat m_3 = -3.706$ en
$\hat m_4 = \qty{-0.931}{ms.km^{-1}}$.

\textbf{18.} $\delta v/v_0 = -v_0\,\delta s = -8.0\times10^{-3}\,\hat m$, in procent
$-0.8\,\hat m$: cel 1 $+0.07$, cel 2 $-1.48$, cel 3 $+2.96$, cel 4
$\qty{+0.75}{\percent}$.

\textbf{19.} Elk diagonaalelement van $\mathsf N^{-1}$ is $11/80$: elke cel heeft
de standaardfout $\sqrt{0.1375} = \qty{0.371}{ms.km^{-1}}$, dat is
$\qty{0.30}{\percent}$. De correlatie van cel 1 en 2 is $-4/11 = -0.36$:
ze delen straal A, en een langzaamheid die in de ene cel wordt gestopt, kan
tegen de andere worden verrekend. Cel 1 en 4, die alleen de minder nauwkeurige
diagonaal E delen, zijn zwak en positief gecorreleerd, $+1/11 = +0.09$.

\textbf{20.} De voorspelde residuen zijn $0.176$, $-0.464$, $-0.379$, $0.091$,
$-0.143$ en $\qty{-0.263}{s}$; de residuen na de fit zijn $-0.016$,
$-0.016$, $0.019$, $0.019$, $-0.007$ en $\qty{-0.007}{s}$, en
$\chi^2 = 0.50$ bij $6 - 4 = 2$ vrijheidsgraden: in overeenstemming met de
opgegeven fouten, eerder aan de kleine kant.

\textbf{21.} De eigenvectoren zijn die van $\mathsf N$ en de eigenwaarden
$\lambda/(\lambda + 4)$: $20/24 = 5/6$, $4/8 = 1/2$, $8/12 = 2/3$ tweemaal. Net
als bij vraag 16 geldt $R_{11} = \tfrac14(\tfrac56 + \tfrac12 + \tfrac23 + \tfrac23) =
\tfrac23$; $R_{12} = \tfrac14(\tfrac56 - \tfrac12 + \tfrac23 - \tfrac23) =
\tfrac1{12}$; $R_{13} = \tfrac14(\tfrac56 - \tfrac12 - \tfrac23 + \tfrac23) =
\tfrac1{12}$; $R_{14} = \tfrac14(\tfrac56 + \tfrac12 - \tfrac23 - \tfrac23) = 0$.

\textbf{22.} Projecteer $\mathbf b$ op de eigenvectoren: $\mathbf b\cdot\mathbf u_k/4
= -14.37$, $0.849$, $12.80$ en $-9.40$ voor de eigenwaarden 20, 4, 8 en 8.
De gedempte schatting deelt ze door $\lambda + 4$ in plaats van door $\lambda$; voor
cel 3 (componenten $1, -1, -1, 1$),
$m_3 = -14.37/24 - 0.849/8 - 12.80/12 - 9.40/12 = \qty{-2.555}{ms.km^{-1}}$,
een snelheidsanomalie van $\qty{+2.04}{\percent}$. Demping trekt elke component
naar nul, het sterkst de slecht bepaalde.

\textbf{23.} Het ongedempte model, dat de data binnen hun fouten past
($\chi^2 = 0.50$ bij 2 vrijheidsgraden), terwijl het gedempte model ze slecht
misfit ($15.8$). Demping is alleen nodig wanneer $\mathsf N$ singulier of bijna
singulier is; hier is de kleinste eigenwaarde, 4, niet klein, de schatting is
stabiel en de fouten zijn matig.

\textbf{24.} Rij 3 is $(1/12, 0, 2/3, 1/12)$: de gedempte schatting van cel 3 is
twee derde van de werkelijke waarde plus een twaalfde van cel 1 en 4, de buren
langs straal C en straal B, en niets van cel 2.

\textbf{25.} Cel 3, de onderste cel links, is ongeveer
\qty{3}{\percent} snel; cel 2 is langzaam. Een koude, dichte slab van oceanische
lithosfeer die onder het continent zinkt, is sneller dan de omringende mantel,
dus de snelle cel, diep en aan de oceaankant van K, is de slab; de langzame
cel erboven aan de andere kant ligt waar de mantelwig heet is.

\textbf{26.} Uit vragen 18 en 19 volgt dat de P-snelheidsanomalie van de slabcel
\textbf{\qty[text-series-to-math=true, uncertainty-mode=separate]{3.0 \pm 0.3}{\percent}} is.
\end{solution}
